\documentclass[aspectratio=169]{beamer}
\usetheme{Madrid}
\usecolortheme{default}
\setbeamertemplate{navigation symbols}{}

\usepackage[utf8]{inputenc}
\usepackage[T5]{fontenc}
\usepackage[vietnamese]{babel}
\usepackage{amsmath,amssymb}
\usepackage{booktabs}
\usepackage{graphicx}
\usepackage{hyperref}
\usepackage{multirow}

\title[Tổng quan hệ thống truy xuất tài liệu]{Tổng quan hệ thống truy xuất tài liệu khoa học}
\author{DALN Project}
\date{\today}

\begin{document}

\begin{frame}
  \titlepage
\end{frame}

\begin{frame}{1. Bài toán}
  \begin{block}{Vấn đề}
    Với một kho dữ liệu khoa học lớn (hơn 13.700 tài liệu), người dùng cần tìm các bài báo phù hợp với truy vấn chuyên ngành như:
    \begin{itemize}
      \item chất liệu thay thế, vật liệu mới, công nghệ pin,\
      \item truy vấn hỗn hợp nhiều thuộc tính,\
      \item kết hợp tiêu chí từ khóa, nguồn, năm, và tính phù hợp ngữ nghĩa.
    \end{itemize}
  \end{block}

  \vspace{0.3cm}
  \begin{block}{Mục tiêu}
    Xây dựng hệ thống truy xuất thông tin để:
    \begin{itemize}
      \item tìm tài liệu chính xác hơn,\
      \item giữ tốc độ phản hồi hợp lý,\
      \item hỗ trợ lọc theo nguồn, năm, và từ khóa bắt buộc,\
      \item trả về Top-$k$ tài liệu kèm đoạn trích và giải thích.
    \end{itemize}
  \end{block}
\end{frame}

\begin{frame}{2. Dữ liệu và tiền xử lý}
  \begin{block}{Dữ liệu đầu vào}
    Mỗi tài liệu chứa các trường quan trọng: title, abstract, source, year, identifier, provenance.
  \end{block}

  \vspace{0.2cm}
  \begin{block}{Tiền xử lý}
    \begin{enumerate}
      \item Lọc và chuẩn hóa văn bản: chữ thường, bỏ ký tự thừa, xử lý từ đồng nghĩa/alias.\
      \item Giữ các trường quan trọng: title + abstract để làm cơ sở truy xuất.\
      \item Tách token và chuẩn hóa chuỗi công thức/thuật ngữ chuyên ngành.\
      \item Tạo biểu diễn lexical và dense cho mỗi tài liệu.\
    \end{enumerate}
  \end{block}

  \vspace{0.2cm}
  \begin{alertblock}{Lưu ý thực tế}
    Tiền xử lý không chỉ là lowercase; cần bảo toàn các thuật ngữ kỹ thuật, tên chất liệu, thủ pháp và chuỗi có tính đặc thù.
  \end{alertblock}
\end{frame}

\begin{frame}{3. Mô hình 1: BM25 (truy xuất từ vựng)}
  \begin{block}{Định nghĩa}
    BM25 là mô hình truy xuất theo từ khóa, dựa trên tần suất từ trong tài liệu và độ hiếm của từ trong corpus.
  \end{block}

  \begin{equation}
    \mathrm{BM25}(q,d) = \sum_{t \in q \cap d} \mathrm{IDF}(t) \cdot \frac{f(t,d)(k_1+1)}{f(t,d)+k_1\left(1-b+b\frac{|d|}{\mathrm{avgdl}}\right)}
  \end{equation}

  \begin{itemize}
    \item $f(t,d)$: tần suất từ $t$ trong tài liệu $d$\
    \item $|d|$: độ dài tài liệu\
    \item $\mathrm{avgdl}$: độ dài trung bình của corpus\
    \item $k_1, b$: tham số điều chỉnh độ nhạy và độ dài tài liệu.
  \end{itemize}

  \vspace{0.2cm}
  \begin{block}{Ý nghĩa}
    BM25 mạnh ở các truy vấn có từ khóa rõ ràng, đặc biệt khi tài liệu chứa các thuật ngữ chuyên ngành xuất hiện nhiều lần.
  \end{block}
\end{frame}

\begin{frame}{4. Mô hình 2: Dense retrieval bằng embedding}
  \begin{block}{Định nghĩa}
    Dense retrieval biễu diễn cả truy vấn và tài liệu thành vector trong không gian ngữ nghĩa tương tự.
  \end{block}

  \begin{equation}
    \mathrm{sim}(q,d)=\cos(\mathbf{e}_q,\mathbf{e}_d)=\frac{\mathbf{e}_q^\top \mathbf{e}_d}{\|\mathbf{e}_q\|\|\mathbf{e}_d\|}
  \end{equation}

  \begin{itemize}
    \item $\mathbf{e}_q$: embedding của truy vấn\
    \item $\mathbf{e}_d$: embedding của tài liệu\
    \item BGE pretrained là mô hình embedding chủ đạo dùng để đo độ tương đồng ngữ nghĩa.\
  \end{itemize}

  \vspace{0.2cm}
  \begin{block}{Ưu điểm}
    Dense retrieval tốt hơn ở trường hợp truy vấn có ý nghĩa tương đương nhưng không trùng đúng từ khóa.
  \end{block}
\end{frame}

\begin{frame}{5. Hybrid retrieval và reranker}
  \begin{block}{Hybrid search}
    Kết hợp đồng thời hai tín hiệu:
    \begin{itemize}
      \item BM25: mạnh về khớp từ khóa\
      \item Dense: mạnh về ngữ nghĩa\
    \end{itemize}
  \end{block}

  \begin{equation}
    \mathrm{RRF}(d) = \sum_{r \in \{\mathrm{BM25},\mathrm{Dense}\}} \frac{1}{k + \mathrm{rank}_r(d)}
  \end{equation}

  \begin{block}{Reranker}
    Sau khi lấy Top-$N$ ứng viên, dùng mô hình rerank để sắp xếp lại theo mức độ liên quan thực sự.
    \begin{equation}
      s_{rerank}(q,d)=f_\theta([q;d])
    \end{equation}
    với $f_\theta$ là mô hình cross-encoder (ví dụ MiniLM reranker).
  \end{block}
\end{frame}

\begin{frame}{6. Quy trình hệ thống}
  \begin{enumerate}
    \item \\textbf{Thu thập và làm sạch dữ liệu}: loại bỏ văn bản không hợp lệ, giữ title + abstract.\
    \item \\textbf{Tiền xử lý văn bản}: chuẩn hóa, token hóa, giữ tên thuật ngữ chuyên ngành.\
    \item \\textbf{Tạo chỉ mục}: xây dựng index BM25 và embedding cho toàn bộ corpus.\
    \item \\textbf{Truy vấn ban đầu}: lấy Top-$N$ tài liệu từ BM25 và dense.\
    \item \\textbf{Hội tụ và sắp xếp}: hợp nhất bằng hybrid fusion, sau đó rerank.\
    \item \\textbf{Lọc bổ sung}: nguồn, năm, từ khóa bắt buộc, các ràng buộc chuyên môn.\
    \item \\textbf{Trả kết quả}: Top-$k$ tài liệu, đoạn trích, và dẫn chứng nội dung.
  \end{enumerate}
\end{frame}

\begin{frame}{7. Mô hình hiệu quả: nhìn từ góc độ toán học}
  \begin{block}{Tổng thể}
    Hệ thống tối ưu một hàm mục tiêu xấp xỉ độ phù hợp giữa truy vấn $q$ và tài liệu $d$:
    \begin{equation}
      \hat{d} = \arg\max_{d \in \mathcal{D}} \; S(q,d)
    \end{equation}
    với
    \begin{equation}
      S(q,d) = \alpha \cdot \mathrm{BM25}(q,d) + \beta \cdot \mathrm{sim}(q,d) + \gamma \cdot \mathrm{Rerank}(q,d)
    \end{equation}
  \end{block}

  \begin{itemize}
    \item $\alpha, \beta, \gamma$: trọng số phức hợp\
    \item Mục tiêu là nâng độ phù hợp tổng hợp giữa khớp từ khóa và ngữ nghĩa.\
  \end{itemize}
\end{frame}

\begin{frame}{8. Kết quả mong đợi / đầu ra hệ thống}
  \begin{block}{Kết quả đầu ra}
    Hệ thống trả về:
    \begin{itemize}
      \item Top-$k$ tài liệu liên quan nhất,\
      \item tiêu đề + abstract,\
      \item năm xuất bản và nguồn,\
      \item đoạn trích phù hợp,\
      \item điểm số tương đồng và lý do lọt top.\
    \end{itemize}
  \end{block}

  \vspace{0.3cm}
  \begin{block}{Giá trị thực tế}
    Hệ thống giúp giảm thời gian tìm kiếm, tăng khả năng tái sử dụng dữ liệu khoa học và hỗ trợ khám phá topic mới trong corpus lớn.
  \end{block}
\end{frame}

\begin{frame}{9. Kết luận}
  \begin{block}{Tóm tắt}
    Mô hình truy xuất tài liệu hiệu quả không chỉ dựa trên từ khóa, mà cần kết hợp:
    \begin{itemize}
      \item tiền xử lý văn bản hợp lý,\
      \item biểu diễn tần suất và ngữ nghĩa,\
      \item hợp nhất nhiều tín hiệu bằng hybrid retrieval,\
      \item reranking để nâng độ chính xác cuối cùng.
    \end{itemize}
  \end{block}

  \vspace{0.3cm}
  \begin{center}
    \Large \textbf{Hệ thống truy xuất tài liệu khoa học là sự kết hợp giữa lexical search, semantic search và reranking.}
  \end{center}
\end{frame}

\end{document}
